\chapter{泡沫解剖}
\label{ch:bubble}

\section{七个泡沫信号}

\subsection{估值与收入的数量级脱节}

Figure AI 以 395 亿美元估值融资，同期\textbf{经常性收入为零}\cite{figure2025seriesc}。银河通用估值 211 亿元，营收推测为千万级——P/S 倍数可能超过 200 倍\cite{galbot2026aplus}。高盛统计：全球约 60\% 的机器人企业估值超出营收 100 倍\cite{gs2025roboticsoutlook}。

更极端的案例是宇树科技：按 420 亿元 IPO 估值、工业场景收入约 3,200 万元计算，真正的「生产力收入」P/S 倍数约为 \textbf{1,300 倍}\cite{unitree2026ipo}。

\subsection{产能规划过载}

仅宇树一家的远期规划产能（7.5 万台人形/年）就已超过 IDC 预测的 2030 年全中国需求（6 万台）\cite{idc2025humanoid,unitree2026ipo}。如果将所有头部企业的产能规划加总，是权威需求预测的 2--3 倍。

\subsection{融资速度异常}

智平方（X Square）1 年内完成 12 轮融资，被称为「全球融资节奏最快的具身智能公司」\cite{xsquare2026}。Manifold AI 成立 10 个月即获近 2 亿元 Pre-A 轮。这种速度更多反映了 VC 之间的 FOMO，而非企业基本面。

\subsection{关联交易疑云}

《财经》杂志报道\cite{caijing2026robots}：很多公司 2025 年收入数千万元，但 2026 年仅两个月就收到上亿订单，「很多是关联交易」。千寻智能创始人韩峰涛直言业内存在「为了融资骗投资人，拿订单骗，或拿开源模型套壳骗」的现象\cite{qianxun2026fraud}。

\subsection{150+ 企业同质化竞争}

中国已有超过 150 家人形机器人企业，2025 年发布 330+ 款机型\cite{miit2026jan}。30+ 家排队 IPO。发改委李超在 2025 年 11 月首次对行业提示风险：「防范重复度高的产品扎堆上市」\cite{ndrc2025warning}。

\subsection{政策驱动的地方互卷}

\begin{table}[htbp]
\centering
\caption{Provincial embodied AI policy competition (targets for 2027).}
\label{tab:policies}
\begin{tabularx}{\textwidth}{lXr}
\toprule
\rowcolor{figLightGray}
\textbf{City} & \textbf{Target} & \textbf{Fund Scale} \\
\midrule
Beijing & 100+ key tech breakthroughs, 10K robots deployed, 100B cluster & National-level \\
Shanghai & 20+ core tech breakthroughs, 50B core industry & Municipal \\
Shenzhen & 10+ 10B-valued companies, 100B industry scale & 20B dedicated fund \\
Chengdu & 50B industry scale & Provincial \\
Suzhou & 200B total robotics industry & 10B fund (Wuzhong) \\
\bottomrule
\end{tabularx}
\end{table}

几乎所有一线、二线城市都在建设机器人产业园、出台专项扶持。北京一个城市三个区各自规划了机器人产业园。这与新能源汽车早期各地蜂拥而上的态势高度类似。

\subsection{分析师预测差异 1,000 倍}

\begin{table}[htbp]
\centering
\caption{Market size predictions --- the 1,000$\times$ divergence.}
\label{tab:predictions}
\begin{tabular}{llr}
\toprule
\rowcolor{figLightGray}
\textbf{Source} & \textbf{Prediction} & \textbf{Timeframe} \\
\midrule
MarketsandMarkets & \$18B & 2030 \\
Goldman Sachs & \$38B & 2035 \\
Fortune Business Insights & \$165B & 2030 \\
Morgan Stanley (bull case) & \$5T & 2050 \\
ARK Invest & \$24T & 2050 \\
\bottomrule
\end{tabular}
\end{table}

从 180 亿到 24 万亿美元，预测之间的差异超过 \textbf{1,000 倍}\cite{gs2024humanoid,morgan2025humanoid100,ark2025humanoid}。当分析师自己都无法在一个数量级内达成共识时，所有预测本质上都是猜测。

\begin{whybox}[Goldman Sachs 的 6 倍上调]
Goldman Sachs 在 2022 年最初预测人形机器人市场为 60 亿美元/10--15 年。到 2024 年，它将预测\textbf{上调 6 倍}至 380 亿美元/2035 年\cite{gs2024humanoid}。两年内基本面并未发生 6 倍变化——发生变化的是叙事的热度。这种行为更接近跟风（momentum），而非独立分析。
\end{whybox}

\section{历史类比}

\subsection{最接近的模板：自动驾驶}

\begin{table}[htbp]
\centering
\caption{Humanoid robots vs. autonomous driving: pattern matching.}
\label{tab:adparallel}
\begin{tabularx}{\textwidth}{lXX}
\toprule
\rowcolor{figLightGray}
\textbf{Dimension} & \textbf{Autonomous Driving (2016--2024)} & \textbf{Humanoid Robots (2024--2026)} \\
\midrule
Total invested & >\$100B & >\$18B (and accelerating) \\
Core promise & ``Full autonomy by 2020'' & ``Million units by 2027'' \\
Celebrity catalyst & Musk (2015: ``FSD by 2018'') & Musk (2021: ``Optimus will change everything'') \\
Reality check & Cruise collapsed 2023; Waymo still limited & Musk admits ``no useful work'' 2025 Q4 \\
Survivors & Waymo (Google), maybe 2--3 others & TBD \\
\bottomrule
\end{tabularx}
\end{table}

自动驾驶累计投入超过 1,000 亿美元\cite{bloomberg2022av}。Musk 在 2015 年承诺 2018 年全自动驾驶。Lyft 联合创始人宣称 2021 年大多数出行将是自动的。Cruise（GM 每年烧 20 亿美元）在 2023 年崩溃。到 2026 年，Waymo 每周完成 45 万次付费出行，估值 1,260 亿美元——但技术有效和商业模式有效是\textbf{完全不同的事情}\cite{waymo2026stats}。

\begin{corebox}[自动驾驶的教训]
技术确实在进步，但比所有预测慢 5--10 年。大多数公司死亡或被收购。最终只有少数 1--3 个赢家存活。这很可能也是人形机器人的轨迹。
\end{corebox}

\subsection{其他参照}

\textbf{元宇宙}（2021--2022）。Meta 改名，投入超 360 亿美元。实际用户远低于预期。两年内从「下一个互联网」变成笑柄。

\textbf{区别}：人形机器人有一个元宇宙缺乏的关键要素——\textbf{真实的劳动力短缺}。全球老龄化和制造业用工荒是不可逆的结构性趋势。这意味着需求端是真实的，问题在供给端。

\section{诚实的声音}

\textbf{Rodney Brooks}（iRobot 联创，MIT）\cite{brooks2025dexterity}：「人形机器人泡沫注定破裂。Optimus 和 Figure 是纯粹的幻想思维。当前方法缺乏触觉反馈——你不能从视频中学会触觉。」

\textbf{朱啸虎}（金沙江创投）\cite{zhuxiaohu2025exit}：2025 年 3 月公开批量退出人形机器人公司，质疑「商业化客户是 CEO 想象出来的」。

\textbf{Gartner}\cite{gartner2026humanoid}：到 2028 年，不到 100 家公司会超越概念验证阶段。不到 20 家会在供应链和制造中实际部署。部署仅限于严格控制的环境。

\textbf{华映资本刘天杰}\cite{huaying2026bubble}：「如果拿一级市场关注度、资金供给去对比商业结果、出货量，人形显然在泡沫阶段，这个不用回避。」

\textbf{发改委}（2025.11）\cite{ndrc2025warning}：首次对行业提示风险——「速度」与「泡沫」需要平衡，要防范重复度高的产品扎堆。这是中国监管层对具身智能发出的最强降温信号。
